\section{Conclusion}
This paper presented a comprehensive study on the design and implementation of a robotic arm for lab testing and education. It discussed the hardware components,
control algorithms, and the integration of sensors for feedback. The results demonstrate the effectiveness of the proposed system in performing complex tasks with high precision
and reliability.

\subsection{Contributions}
In summary, the hardware improvements made to the robotic arm include:

\begin{itemize}
    \item Replacement of 3D-printed axes with aluminum axes for increased durability.
    \item Introduction of ball bearings to support the axes and reduce friction.
    \item Implementation of a clamping mechanism for better modularity and easier assembly.
    \item Modification of the yaw joint to include a 8:1 planetary gearbox to increase available torque.
    \item Enhancement of the tensioning mechanism for the elbow joint to improve reliability.
    \item Integration of a gripper with manual control, with potential for future software control.
\end{itemize}

The software contributions include:

\begin{itemize}
    \item Development of a torque measurement method using the current sensing capabilities of the servo motors and the results from extensive experiments.
    \item Extension of the existing robot control framework to support current control mode and torque commands.
    \item Implementation of various torque-based control strategies, namely joint space PD-control, Cartesian space PD-control, and end-effector force control.
\end{itemize}

\subsection{Limitations}
Despite the significant improvements made, there are still some limitations to the current design.
Mainly, the stiffness of the robotic arm is still limited by the use of 3D-printed components in some areas. Therefore, the maximum load the arm can handle is still relatively low
if the precision of the end-effector position is to be maintained.

Additionally, the torque measurement method relies on the measured current of the servo motors, which leads to inaccuracies due to unmodeled gearbox friction and other factors.
To be precise, the torque measuremt is ambiguous in the static case, as the same current can correspond to different torques depending on the previous movement and friction.
This limitation could affect the performance of future torque-based control strategies relying on this measurement.

\subsection{Future Work}
Future work could focus on further improving the hardware design to increase the stiffness and load capacity of the robotic arm.
Concrete ideas for further hardware improvements can be found in the Chapter~\nameref{ch:hardware}.

On the software side, future work could focus on refining the torque measurement method to improve accuracy. This could involve developing a more sophisticated model of the motor and
gearbox dynamics or including additional sensors for direct torque measurement.

Furthermore, advanced control strategies using torque sensor data could be explored, such as model-based control or learning-based approaches, to enhance the performance and
adaptability of the robotic arm in various tasks. Specifically, implementing advanced impedance control or admittance control could be beneficial for tasks involving interaction with
the environment.

Lastly, advanced path planning and motion planning algorithms could be integrated into the control framework to enable more complex and dynamic tasks for the robotic arm.